
%===================================================================================================
%===================================================================================================
\section{Introduction}
\label{sec:introduction}
%===================================================================================================
%===================================================================================================

In recent years, deep learning has seen significant growth, driven by several methodologies which enable the training of deep neural networks (DNNs) in a scalable way and by development of more powerful hardwares. It is observed that increased capacity of DNN effectively has improved the performance. For example, AmoebaNet-B \cite{real2019regularized} scaled with GPipe \cite{huang2019gpipe} has 557 million parameters and has achieved top-1 accuracy 84.4\% which was state-of-the-arts result at the time, and GPT-2 \cite{radford2019language} is a Transformer-based \cite{vaswani2017attention} language model which has 1.5 billion parameters (see Figure 1 of \cite{huang2019gpipe} for the effect of model scaling). However, training such a massive model is very resource intensive. One can mitigate this issue by reducing the size of the model without losing the performance by pruning the model \cite{han2015learning, alvarez2016learning}, designing more efficient architectures \cite{howard2017mobilenets, tan2019efficientnet}, architecture search under resource constraints \cite{cai2018proxylessnas}, and many more. 

We may wonder a rather direct approach is possible: \emph{can we train a massive model fast enough, given a large pool of devices?} One obstacle is that common optimization techniques to train a neural network are sequential in nature. Those algorithms repeatedly compute the gradient of the loss with respect to the given mini-batch at a time and update the model parameters using the gradient. With abundant computational resource, data parallelism \cite{krizhevsky2012imagenet} is commonly used to speed up the overall optimization procedure by dividing the mini-batch into micro-batches and delegating per micro-batch computation to available devices. With careful hyperparameter tuning, this effectively reduce the training time up to a certain size of mini-batch which may depend on model, optimization algorithm, and data \cite{goyal2017accurate, shallue2018measuring}. One drawback of data-parallel training is that devices hold their own version of network for executing the subdivided task, and network parameters must be synchronized after each parameter update. This may induce heavy communication load when there are lots of parameters to synchronize.

Note that data parallelism is not applicable when the model is so big that it is impossible to compute gradient even when a single data point is fed into the network. Model parallelism \cite{dean2012large} is a method for training such a massive model, which partitions the model into several pieces and places them on different devices. Each device only computes a small part of the model, and updates only the parameters in that part. However, model parallelism suffers from its underutilization behavior. Since most neural networks consist of sequence of layers, the device holding the later part of the model must wait until computation in devices holding earlier parts of the model. 

Another possible solution is to use gradient checkpointing \cite{chen2016training} which saves memory by only storing the subset of activation maps and re-computing the discarded activation maps when necessary. Obviously, this requires certain part of the model be computed twice and overall training time would be increased. 

It is benefitting to combine different types of parallelization strategies \cite{krizhevsky2014one, pmlr-v80-jia18a, shazeer2018mesh, huo2018decoupled, harlap2018pipedream, huang2019gpipe, guan2019xpipe}, and recent lines of research questions how to find an optimal strategy \cite{jia2018beyond, mirhoseini2017device, mirhoseini2018a, zhou2019gdp}. Among them, pipeline parallelism a way to accelerate neural network training by combining model parallelism with data pipelining, either in synchronous way as in GPipe \cite{huang2019gpipe} or in asynchronous way as in  \cite{huo2018decoupled}, PipeDream \cite{harlap2018pipedream}, and XPipe \cite{guan2019xpipe}. We remark that gradient checkpointing (also called re-materialization) is further combined in GPipe to allow training even bigger models.

In this paper, we design and implement {\torchgpipe}, a ready-to-use library for GPipe in PyTorch \cite{paszke2017automatic}. In particular, we develop a set of design components for optimized pipeline-parallel computations in PyTorch's \emph{define-by-run} and \emph{eager execution} environment. We show that each component is necessary to fully benefit from pipeline parallelism in such environment, and demonstrate the efficiency of {\torchgpipe} by conducting the speed and memory benchmarks on AmoebaNet-D \cite{real2019regularized} and U-Net \cite{RonnebergerFB15} when trained with the library.

The rest of the paper is organized as follows.
In \autoref{sec:pipeline-parallelism}, we discuss how the forward and backward passes can be decomposed into subtasks (under certain assumptions), describe the device placement strategy of micro-batch pipeline parallelism, and demonstrate what the desired order of execution per device is. In \autoref{sec:torchgpipe}, we discuss complications for achieving the optimal timeline of pipeline parallelism in PyTorch and explain how {\torchgpipe} resolves them. Additionally, we relax the assumption that the model is sequentially composed, and provide a way for expressing models with long skip connections so that pipeline parallelism still applies without giving up the efficiency. Then, we demonstrate that the optimization components suggested in the paper are essential for the performance, and evaluate the performance of the proposed library in \autoref{sec:experiments}.
